\begin{table}[h]
\centering
\caption{Sample sizes inside the CMA after elevation matching.}
\label{tab:samples}
\begin{tabular}{llcc}
\toprule
Test & Class & Points & Mean elevation (m) \\
\midrule
Within-2015 spatial CV & Flood (NRSC code 13) & 186 & 10.6 \\
 & Non-flood & 186 & 10.6 \\
Transfer to 2005 & Flood (IRS 2005) & 149 & 14.3 \\
 & Non-flood & 149 & 14.4 \\
\bottomrule
\end{tabular}
\end{table}

\begin{table*}[h]
\centering
\caption{Predictor summaries inside the CMA, median (interquartile range), by event and class.}
\label{tab:summary}
\begin{tabular}{llrlllll}
\toprule
Event & Class & n & Elevation (m) & Slope ($^\circ$) & TWI & Distance to water (m) & Built-up (\%) \\
\midrule
2015 CMA & flood & 186 & 10.2 (5.1--15.9) & 2.0 (1.3--3.1) & 9.7 (9.0--11.5) & 1604.2 (750.4--2859.0) & 65 \\
2015 CMA & non-flood & 186 & 10.2 (5.0--15.6) & 1.1 (0.5--2.0) & 10.0 (9.1--11.0) & 1084.4 (521.3--1815.2) & 36 \\
2005 CMA & flood & 149 & 14.2 (7.0--21.2) & 1.0 (0.5--1.9) & 9.9 (9.2--11.4) & 1374.1 (849.6--2106.4) & 42 \\
2005 CMA & non-flood & 149 & 14.1 (6.9--21.6) & 1.4 (0.7--2.4) & 9.9 (9.0--11.2) & 1405.2 (692.6--2308.2) & 41 \\
\bottomrule
\end{tabular}
\end{table*}

\begin{table}[h]
\centering
\caption{Within-2015 spatial CV ROC-AUC inside the CMA (186 per class), with 95\% bootstrap intervals.}
\label{tab:within}
\begin{tabular}{lcc}
\toprule
Predictors & Logistic regression & Random forest \\
\midrule
elevation only & 0.39 (0.29--0.49) & 0.30 (0.23--0.37) \\
built-up only & 0.59 (0.51--0.67) & 0.58 (0.50--0.67) \\
terrain, no elevation & 0.72 (0.64--0.79) & 0.74 (0.68--0.80) \\
terrain + built-up & 0.74 (0.66--0.81) & 0.75 (0.68--0.82) \\
\bottomrule
\end{tabular}
\end{table}

\begin{table}[h]
\centering
\caption{Transfer from 2015 to 2005 inside the CMA (149 per class). ROC-AUC with 95\% bootstrap intervals.}
\label{tab:transfer}
\begin{tabular}{lcc}
\toprule
Predictors & Logistic regression & Random forest \\
\midrule
elevation only & 0.50 (0.43--0.57) & 0.52 (0.46--0.58) \\
built-up only & 0.51 (0.45--0.56) & 0.51 (0.45--0.56) \\
terrain, no elevation & 0.43 (0.37--0.49) & 0.40 (0.34--0.47) \\
terrain + built-up & 0.45 (0.38--0.51) & 0.42 (0.35--0.48) \\
\bottomrule
\end{tabular}
\end{table}

\begin{table}[h]
\centering
\caption{Label comparison across the study box: hit rates at reported hotspots and stagnation points, against random background points.}
\label{tab:labels}
\begin{tabular}{lcc}
\toprule
Comparison & Reports & Random background \\
\midrule
GCC hotspots on built-up cells & 90\% & 7\% \\
GCC hotspots within 300\,m of NRSC code 13 & 72\% & 28\% \\
GCC hotspots within 300\,m of NRSC code 1 & 25\% & 27\% \\
GCC hotspots on UNOSAT 24 Nov flood cells & 0\% & 23\% \\
Stagnation points within 300\,m of NRSC code 13 & 65\% & 2\% \\
\bottomrule
\end{tabular}
\end{table}
